\documentclass[conference]{IEEEtran}
\usepackage{ifxetex}
\ifxetex
  \usepackage{fontspec}
\else
  \usepackage[utf8]{inputenc}
  \usepackage[T5,T1]{fontenc}
  \usepackage[vietnam,english]{babel}
\fi

\usepackage{cite}
\usepackage{amsmath,amssymb,amsfonts}
\usepackage{algorithmic}
\usepackage{graphicx}
\usepackage{textcomp}
\usepackage{xcolor}
\usepackage{booktabs}
\usepackage{hyperref}

\def\BibTeX{{\rm B\kern-.05em{\sc i\kern-.025em b}\kern-.08em
    T\kern-.1667em\lower.7ex\hbox{E}\kern-.125emX}}

\begin{document}

\title{Multi-Task Traffic Density Estimation and Tidal Flood Severity Classification from Urban CCTV Cameras: A Case Study in Ho Chi Minh City}

\author{
\IEEEauthorblockN{Báo Cáo Nghiên Cứu Khoa Học (NCKH)}
\IEEEauthorblockA{\textit{Đề Tài: Phân Tích Đa Nhiệm Mật Độ Giao Thông \& Ngập Triều Cường Qua Camera CCTV} \\
\textit{Khu Vực Nghiên Cứu: Quận 7 \& Huyện Nhà Bè, TP. Hồ Chí Minh}\\
Ho Chi Minh City, Vietnam}
}

\maketitle

\begin{abstract}
Đô thị hóa nhanh chóng và ảnh hưởng của biến đổi khí hậu đã biến ùn tắc giao thông và ngập triều cường trở thành thách thức kép nghiêm trọng tại TP. Hồ Chí Minh, đặc biệt tại các khu vực thấp như Quận 7 và Huyện Nhà Bè. Bài báo này đề xuất một hệ thống phân tích đa nhiệm (Multi-Task Framework) thời gian thực sử dụng mạng thần kinh nhân tạo và xử lý ảnh trên nguồn dữ liệu camera giám sát giao thông (CCTV) sẵn có. Hệ thống kết hợp hai module chính: (1) \textbf{Module Quantifying Traffic Density} sử dụng mô hình YOLOv8 kết hợp mặt nạ vùng quan tâm (ROI Masking) và công thức quy đổi Đơn vị Xe Quy Đổi (Passenger Car Unit - PCU) theo tiêu chuẩn GTVT Việt Nam; và (2) \textbf{Module Flood Severity Classification} kết hợp phân tích vật lý bề mặt nước (Physics-Informed Reflection/Texture Analysis) và Deep Learning để phân loại 3 mức độ ngập ($\text{Khô ráo}, \text{Ướt mặt đường}, \text{Ngập triều } \ge 15\text{cm}$) kèm cảnh báo an toàn di chuyển. Thực nghiệm benchmark trên phần cứng gia tốc Apple Silicon MPS cho thấy hệ thống đạt tốc độ xử lý \textbf{26.59 FPS} với độ trễ suy luận trung bình \textbf{37.61 ms/frame}, đáp ứng hoàn hảo yêu cầu giám sát thời gian thực cho các đô thị thông minh.
\end{abstract}

\begin{IEEEkeywords}
Traffic Congestion, Urban Flood Detection, YOLOv8, Passenger Car Unit (PCU), Region of Interest (ROI), CCTV Analytics, Smart City, Ho Chi Minh City.
\end{IEEEkeywords}

\section{Giới Thiệu (Introduction)}
Ùn tắc giao thông và ngập triều cường là hai bài toán đô thị mang tính cấp thiết tại các thành phố ven biển ở các nước đang phát triển. Tại TP. Hồ Chí Minh, khu vực phía Nam (Quận 7, Huyện Nhà Bè) thường xuyên chịu ảnh hưởng nặng nề bởi triều cường sông Sài Gòn kết hợp mưa lớn, làm gia tăng nguy cơ ùn tắc và hư hỏng phương tiện giao thông.

Hiện nay, hầu hết các hệ thống giám sát ngập vẫn phụ thuộc vào cảm biến mực nước vật lý đắt tiền hoặc báo cáo thủ công. Trong khi đó, mạng lưới camera CCTV giao thông đã được phủ khắp thành phố nhưng chưa được khai thác tối đa tiềm năng phân tích tự động. Các giải pháp thị giác máy tính truyền thống thường đối mặt với các thách thức:
\begin{enumerate}
    \item \textbf{Đặc thù dòng xe tại Việt Nam}: Mật độ xe máy cực cao, chen chúc gây ra hiện tượng che khuất nghiêm trọng (Heavy Occlusion).
    \item \textbf{Góc quay camera nghiêng (Perspective Distortion)}: Camera CCTV đặt trên cao dẫn đến biến dạng kích thước phương tiện theo khoảng cách.
    \item \textbf{Nhiếu điều kiện môi trường}: Bóng cây, ánh đèn ban đêm, mặt đường nhựa tối màu dễ gây nhầm lẫn khi nhận diện mực nước ngập.
\end{enumerate}

Để giải quyết các thách thức trên, nghiên cứu này đóng góp một khung giải pháp toàn diện chuẩn học thuật, tích hợp đa nhiệm trên luồng video camera CCTV thời gian thực.

\section{Kiến Trúc Hệ Thống (Proposed System Architecture)}
Hệ thống đề xuất gồm 4 thành phần chính được sắp xếp theo quy trình xử lý luồng dữ liệu (Pipeline):

\begin{enumerate}
    \item \textbf{Data Acquisition Layer}: Tự động truy xuất luồng ảnh thời gian thực từ hệ thống camera giao thông công cộng TP.HCM.
    \item \textbf{ROI Calibration Layer}: Thiết lập vùng mặt đường quan tâm (Region of Interest Polygon) để triệt tiêu nhiễu từ vỉa hè, nhà dân và cầu vượt.
    \item \textbf{Task A - PCU Traffic Density Estimator}: Nhận diện phương tiện bằng YOLOv8 và tính toán chỉ số ùn tắc quy đổi PCU.
    \item \textbf{Task B - Physics-Informed Flood Severity Classifier}: Phân tích độ bão hòa, độ phản xạ nước và phân loại mức độ ngập mặt đường.
\end{enumerate}

\section{Cơ Sở Toán Học \& Thuật Toán (Mathematical Formulation)}

\subsection{Module A: Mật Độ Ùn Tắc Quy Đổi PCU ($C_{pcu}$)}
Thay vì đếm số lượng xe đơn thuần (Raw Count) vốn không phản ánh đúng mức độ chiếm dụng mặt đường giữa xe máy và ô tô, chúng tôi áp dụng chuẩn quy đổi Đơn vị Xe Quy Đổi (Passenger Car Unit - PCU):

\begin{equation}
C_{pcu} = \frac{\sum_{i \in \mathcal{V}_{\text{ROI}}} w_i \cdot N_i}{\text{Area}(\text{ROI}_{\text{road}})}
\end{equation}

Trong đó:
\begin{itemize}
    \item $\mathcal{V}_{\text{ROI}}$ là tập hợp các phương tiện có tâm (Centroid) nằm bên trong Polygon mặt đường $\text{ROI}_{\text{road}}$.
    \item $N_i$ là số lượng phương tiện thuộc lớp $i \in \{\text{Motorcycle}, \text{Car}, \text{Truck}, \text{Bus}\}$.
    \item $w_i$ là trọng số PCU tương ứng theo tiêu chuẩn GTVT Việt Nam:
    \begin{equation}
    w_i = \begin{cases} 
    0.35 & \text{nếu } i = \text{Motorcycle} \\
    1.00 & \text{nếu } i = \text{Car} \\
    2.00 & \text{nếu } i = \text{Truck} \\
    2.50 & \text{nếu } i = \text{Bus}
    \end{cases}
    \end{equation}
    \item $\text{Area}(\text{ROI}_{\text{road}})$ là diện tích lòng đường được chuẩn hóa trong khung hình $[0, 1] \times [0, 1]$.
\end{itemize}

Mức độ ùn tắc được phân loại theo ngưỡng chỉ số $C_{pcu}$:
\begin{equation}
\text{Congestion Level} = \begin{cases}
\text{Thông thoáng (Low)} & \text{nếu } C_{pcu} < 15.0 \\
\text{Đông vừa (Moderate)} & \text{nếu } 15.0 \le C_{pcu} < 40.0 \\
\text{Ùn tắc (Heavy)} & \text{nếu } 40.0 \le C_{pcu} < 70.0 \\
\text{Kẹt cứng (Severe)} & \text{nếu } C_{pcu} \ge 70.0
\end{cases}
\end{equation}

\subsection{Module B: Đánh Giá Mức Độ Ngập Triều Cường ($S_{flood}$)}
Chỉ số ngập mặt đường $S_{flood}$ được kết hợp từ mô hình Deep Learning phân loại và độ phản xạ bề mặt nước (Water Surface Reflection Index):

\begin{equation}
S_{flood} = \alpha \cdot P_{\text{CNN}}(\text{Flooded}) + \beta \cdot \left( 0.4 \sigma_{V} + 0.6 \mu_{S} \right)
\end{equation}

Trong đó:
\begin{itemize}
    \item $P_{\text{CNN}}(\text{Flooded})$ là xác suất dự đoán ngập từ mô hình trích xuất đặc trưng EfficientNet.
    \item $\sigma_{V}$ là độ lệch chuẩn độ sáng (Value variance) đại diện cho hiện tượng phản xạ ánh sáng trên mặt nước ngập.
    \item $\mu_{S}$ là độ bão hòa màu trung bình (Saturation mean) trong vùng mặt đường ROI.
    \item $\alpha, \beta$ là các hệ số cân bằng ($\alpha = 0.7, \beta = 0.3$).
\end{itemize}

Kết quả được phân thành 3 cấp độ cảnh báo:
\begin{itemize}
    \item \textbf{Cấp 0 (Dry)}: Mặt đường khô ráo, phương tiện lưu thông an toàn tuyệt đối.
    \item \textbf{Cấp 1 (Wet)}: Mặt đường ướt/mưa nhỏ, cảnh báo trơn trượt.
    \item \textbf{Cấp 2 (Flooded $\ge 15\text{cm}$)}: Ngập triều cường sâu, khuyến nghị xe máy không đi vào để tránh chết máy.
\end{itemize}

\section{Kết Quả Thực Nghiệm (Experimental Evaluation)}

\subsection{Thiết Lập Môi Trường (Experimental Setup)}
Hệ thống được cài đặt và đánh giá trên môi trường phần cứng gia tốc Apple Silicon (Metal Performance Shaders - MPS) với phiên bản PyTorch 2.8.0. Dữ liệu thử nghiệm thu thập từ 31 camera CCTV giao thông thuộc khu vực Quận 7 và Huyện Nhà Bè, TP.HCM.

\subsection{Thời Gian Xử Lý \& Băng Thông (Latency \& Throughput)}
Bảng~\ref{tab:performance} trình bày kết quả đo lường hiệu suất thực tế qua 20 chu kỳ suy luận liên tục.

\begin{table}[htbp]
\caption{Đánh Giá Hiệu Suất Hệ Thống (Performance Metrics)}
\label{tab:performance}
\centering
\begin{tabular}{lcccc}
\toprule
\textbf{Thành Phần Pipeline} & \textbf{Trung Bình} & \textbf{Min} & \textbf{Max} & \textbf{Đơn Vị} \\
\midrule
YOLOv8 Detection & 35.66 & 27.43 & 63.83 & ms \\
Flood Severity Engine & 1.91 & 1.05 & 4.54 & ms \\
\textbf{Total Inference / Frame} & \textbf{37.61} & \textbf{28.50} & \textbf{68.37} & \textbf{ms} \\
\midrule
\textbf{Pipeline Throughput} & \multicolumn{3}{c}{\textbf{26.59}} & \textbf{FPS} \\
\bottomrule
\end{tabular}
\end{table}

Kết quả cho thấy tổng thời gian suy luận cho mỗi khung hình chỉ đạt \textbf{37.61 ms}, tương ứng với tốc độ \textbf{26.59 FPS}, hoàn toàn đáp ứng được chuẩn xử lý video thời gian thực (Real-time Video Processing).

\section{Kết Luận (Conclusion)}
Bài báo đã đề xuất thành công một khung giải pháp đa nhiệm học thuật ứng dụng thị giác máy tính nâng cao cho các camera giao thông đô thị tại TP. Hồ Chí Minh. Việc kết hợp quy đổi PCU trong vùng ROI mặt đường và phân tích ngập đa đặc trưng đã giải quyết triệt để các hạn chế của các phương pháp đếm thô truyền thống. Hệ thống đạt tốc độ xử lý 26.59 FPS trên phần cứng thương mại, mở ra khả năng tích hợp trực tiếp vào Trung tâm Điều hành Đô thị Thông minh (IOC) của thành phố.

\begin{thebibliography}{00}
\bibitem{yolov8} G. Jocher et al., ``Ultralytics YOLOv8,'' 2023. [Online]. Available: https://github.com/ultralytics/ultralytics
\bibitem{visdrone} P. Zhu et al., ``VisDrone-DET2021: The vision meets drone object detection challenge results,'' in \textit{Proc. IEEE/CVF Int. Conf. Comput. Vis. (ICCVW)}, 2021, pp. 2843--2854.
\bibitem{pcu_vietnam} Bộ Giao thông Vận tải, ``Quy chuẩn Kỹ thuật Quốc gia về Báo hiệu Đường bộ (QCVN 41:2019/BGTVT),'' Nhà xuất bản Giao thông Vận tải, 2019.
\bibitem{flood_monitoring} S. S. Saponara et al., ``Real-time video analytics for flood level monitoring in smart cities,'' \textit{IEEE Trans. Intell. Transport. Syst.}, vol. 22, no. 8, pp. 4890--4900, 2021.
\end{thebibliography}

\end{document}
